\documentclass[12pt]{article}
\usepackage[a4paper,margin=2.3cm]{geometry}
\usepackage{xcolor}
\usepackage{parskip}
\usepackage{microtype}

\definecolor{navy}{HTML}{1B2A4A}
\definecolor{brick}{HTML}{C0392B}
\definecolor{muted}{HTML}{6B7280}
\definecolor{tint}{HTML}{EEF3FB}

% slide-cue header: #1 = slide name, #2 = running time
\newcommand{\cue}[2]{%
  \par\medskip
  {\color{navy}\rule{\linewidth}{0.6pt}}\par\vspace{1pt}%
  \noindent{\sffamily\bfseries\color{navy}\large SLIDE \textbullet\ #1}%
  \hfill{\sffamily\small\color{muted} #2}\par\smallskip}
% stage direction (not spoken)
\newcommand{\sd}[1]{{\color{muted}\itshape (#1)}\ }
% emphasis word to stress aloud
\newcommand{\key}[1]{\textbf{#1}}

\pagestyle{empty}
\linespread{1.2}

\begin{document}

% ---------- Title block ----------
{\sffamily
 {\color{navy}\Huge\bfseries Cadence}\\[3pt]
 {\color{brick}\large\bfseries Software segment --- presentation script}\\[5pt]
 {\color{muted}\normalsize 4-speaker group talk \textbullet\ your part \textbullet\ target $\approx$ 3 minutes}\par}
\vspace{5pt}
{\color{navy}\rule{\linewidth}{1.2pt}}

% ---------- Reminders strip ----------
\vspace{7pt}
\noindent\colorbox{tint}{%
\begin{minipage}{\dimexpr\linewidth-2\fboxsep}
\sffamily\small\color{navy}\setlength{\parskip}{2pt}
\textbf{Before you start}\\
$\bullet$\ Let \emph{Hardware} hand you ``a stream of samples'' --- \emph{you} own the Freeze Index, gate \& debounce.\\
$\bullet$\ Frame your pipeline as the \emph{hand-built, rule-based} detector (it sets up the ML handoff).\\
$\bullet$\ Running long? The ``Host tooling'' slide is the safe cut. \quad Bold words = stress them.
\end{minipage}}
\vspace{4pt}

% ---------- Script ----------
{\large\raggedright

\cue{from raw motion to a live decision}{0:00}
Thanks. So \textsf{[Speaker~2]} hands the software a stream of accelerometer samples at \key{64 hertz}. My job is everything that turns that raw stream into a freeze decision you can \key{trust} --- both on the device in real time, and on the host afterwards.

\cue{on-device pipeline}{0:25}
Here's what runs on the board itself. Every 2 seconds, I take the last 4-second window --- 256 samples --- and run a short \key{FFT} over it. I sum the power in the freeze band, 3 to 8 hertz, and the power in the walking band, half to 3 hertz, and divide them --- that ratio is the \key{Freeze Index}.

But --- and this is the important part --- a high index on its own is \emph{not} a freeze. The decision logic only says \emph{freeze} when \key{two conditions} are met: there's enough overall movement energy to clear the \key{gate}, \emph{and} the index is over \key{threshold}. Then a \key{two-window debounce} state machine requires that to hold for two readings in a row before it flips the state, on or off. That single piece of logic is what kills the one-off false alarms. The result drives the on-board display live --- and yes, the cue path exists in the code, but it's \key{disabled in firmware} by design.

And all of that --- an FFT and a decision, twice a second --- runs in lightweight integer math on a tiny low-power microcontroller. No phone, no cloud.

\cue{host tooling \& data pipeline}{1:40}
Off the device, the software does two more jobs. The firmware tags \key{every sample with an auto-incrementing phase label} and logs it to flash; we replay it over USB straight into a CSV, so we get \key{ground truth for free}. And a Python pipeline windows that data, builds the confusion matrices and the threshold sweep --- that's what let us tune the threshold \key{honestly}, not by eye. The same decision stream also feeds a \key{live web dashboard} --- which I'll show you right now.

\cue{live demo}{2:10}
\sd{start the dashboard replay} This is a real capture playing back. Watch the Freeze Index trace against the operating line. Through normal walking it stays \key{under the line} --- silent. And here, the genuine freeze --- it jumps above, the state flips to \textcolor{brick}{\textbf{FREEZE}}, the dashboard lights up. Quiet on the walk, clean catch on the freeze --- that's the gate and the debounce, in software, doing exactly their job.

\cue{close}{2:50}
So the whole pipeline is deliberately \key{transparent} --- rule-based, traceable, every decision explainable. That's the philosophy behind Cadence: \emph{a glass-box monitor a clinician can trust beats a black box that just buzzes.} Thank you --- happy to take questions.

}% end large

\end{document}
